\documentclass[11pt]{article}
\usepackage[localtheorems]{raeez-math-template}

\newtheorem{thm}{Theorem}[section]
\newtheorem{prop}[thm]{Proposition}
\newtheorem{lem}[thm]{Lemma}
\newtheorem{cor}[thm]{Corollary}
\theoremstyle{definition}
\newtheorem{defn}[thm]{Definition}
\newtheorem{rem}[thm]{Remark}
\newtheorem{fact}[thm]{Fact}
\theoremstyle{remark}
\newtheorem{attack}[thm]{Attack}
\newtheorem{heal}[thm]{Heal}

\newcommand{\bC}{\mathbb{C}}
\newcommand{\bR}{\mathbb{R}}
\newcommand{\bZ}{\mathbb{Z}}
\newcommand{\cA}{\mathcal{A}}
\newcommand{\cB}{\mathcal{B}}
\newcommand{\cC}{\mathcal{C}}
\newcommand{\cZ}{\mathcal{Z}}
\newcommand{\cF}{\mathcal{F}}
\newcommand{\cE}{\mathcal{E}}
\newcommand{\fg}{\mathfrak{g}}
\newcommand{\fgl}{\mathfrak{gl}}
\newcommand{\Rep}{\mathrm{Rep}}
\newcommand{\Hoch}{\mathrm{HH}}
\newcommand{\Eone}{E_1}
\newcommand{\Etwo}{E_2}
\newcommand{\Ethree}{E_3}
\newcommand{\Einf}{E_\infty}
\newcommand{\Conf}{\mathrm{Conf}}
\newcommand{\Ran}{\mathrm{Ran}}
\newcommand{\HH}{\mathrm{HH}}
\newcommand{\Zder}{\cZ^{\mathrm{der}}_{\mathrm{ch}}}
\newcommand{\SCchtop}{\mathsf{SC}^{\mathrm{ch,top}}}
\newcommand{\Sug}{\mathrm{Sug}}
\newcommand{\barB}{\bar B}
\newcommand{\id}{\mathrm{id}}
\DeclareMathOperator{\res}{res}
\DeclareMathOperator{\av}{av}

\title{The $E_n$ operadic spiral:\\
five arrows, one circle, three volumes.\\
Adversarial heal on \texttt{opus\_08}}
\author{Wave~17 elite-voice adversarial-heal}
\date{2026-04-24}

\begin{document}
\maketitle

\paragraph{Voice.}
Primary: \textsc{Drinfeld} --- unifying principle, associator,
quantum double. Secondary: \textsc{Bezrukavnikov} --- geometric
representation theory, derived category, sheaf-theoretic
Koszul. Tertiary: \textsc{Gaiotto} --- dualities compute.

\paragraph{Target.}
The five-arrow passage
\[
\Ethree^{\mathrm{top}}(\text{bulk})
\xrightarrow{\;\mathrm{res}\;}
\Etwo\text{-chiral}(\text{bdy})
\xrightarrow{\;\barB^{\mathrm{ord}}\;}
\Eone\text{-chiral}(\text{bar/QG})
\xrightarrow{\;\cZ(\Rep(-))\;}
\Etwo\text{-braided}(\text{Drinfeld centre})
\xrightarrow{\;\HH^\bullet_{\mathrm{cat}}\;}
\Ethree^{\mathrm{top}}(\text{derived centre}),
\]
and the claim that $\pi_1(\Conf_2(\bR^3)) = 0$ forces the
Drinfeld-centre arrow (arrow~3 above) to be the unique source of
non-trivial braiding. The corollary is that three volumes are
forced: two cannot close the circle because no two of these three
geometric settings provides the $\Eone$-chiral source; four would
split $\SCchtop$ from $\Ethree$-topologisation along a fake seam.

\section*{Preamble: distinguish $\Etwo\text{-chiral}$ from
$\Eone\text{-chiral}$}

Before the attack cycles the three $\Etwo$ objects that collide in
the diagram must be pinned. The body of the programme uses three
different $\Etwo$'s, all carrying a $z$-variable, and they must
not be conflated.

\begin{defn}[Three $\Etwo$'s; Vol~I working convention]
\label{def:three-e2}
\begin{enumerate}[label=\textup{(\Alph*)}]
\item $\Etwo$-\emph{holomorphic}. An algebra over the holomorphic
 refinement of the little $2$-disks operad on a complex curve~$C$.
 Equivalent data: a factorisation algebra on $\Ran(C)$ with
 $\Einf$-compatible (symmetric) OPE. Examples: Heisenberg,
 Virasoro, affine Kac--Moody at non-critical level. The
 $z$-variable is the complex coordinate on~$C$.
\item $\Etwo$-\emph{chiral}, class \textup{(B)}. An $A_\infty$
 algebra in $\End^{\mathrm{ch}}_\cA$ after restriction of a
 $3$d HT bulk observable to a codimension-$2$ defect curve.
 Same $z$, but the $\Etwo$-structure records holomorphic
 translation along the curve plus the residual
 $\bR^2$-transverse direction collapsed onto the defect. This is
 the object of arrow~2.
\item $\Etwo$-\emph{braided}. A braided monoidal $(\infty,1)$-category,
 output of the Drinfeld centre arrow~3. The $z$-variable here is
 the spectral parameter on evaluation modules,
 $z = u - v \in \bC^\times$, not a worldsheet coordinate.
\end{enumerate}
\end{defn}

The three $\Etwo$'s live in distinct ambient categories:
$\Etwo$-holomorphic in $\mathrm{Alg}_{\Etwo}(\mathrm{ShvFact}^{\mathrm{hol}}(C))$
(a presentable stable $(\infty,1)$-category of algebras);
$\Etwo$-chiral in the same ambient but with the bulk-restriction
structure that records the transverse geometry;
$\Etwo$-braided in
$\mathrm{BrMon}\text{-Cat}(\infty,1)$, a braided
monoidal $(\infty,1)$-category. Arrow~2 outputs (B); arrow~3
outputs (C); the $\Eone$ in the middle is the algebra on which the
representation category (C) is built. Conflating (A) with (B) is
AP163/AP167 (SC is not $\Ethree$; Dunn additivity does not apply
to coloured operads); conflating (B) with (C) is the cached AP at
\texttt{notes/first\_principles\_cache\_comprehensive.md} entry
\textit{E$_2$-algebra vs E$_2$-category}.

\section{Five attack-heal cycles}

\subsection*{Attack~1 --- the $\pi_1$ claim: literal triviality or
structural triviality?}

\begin{attack}[Is $\pi_1(\Conf_2(\bR^3)) = 0$ the precise form of
``trivial braiding for $\Ethree \to \Etwo$?'']
\label{att:pi1}
The operadic
restriction $\Ethree \to \Etwo$ factors through the little-disks
forgetful $\mathrm{Disk}^{\mathrm{or,fr}}_3 \to
\mathrm{Disk}^{\mathrm{or,fr}}_2$. The braiding on an $\Etwo$-algebra
is the action of $\pi_1$ of the framed configuration space
$\Conf_2^{\mathrm{fr}}$. For $n = 2$ one has
$\pi_1(\Conf_2(\bR^2)) = \pi_1(S^1) = \bZ$, the source of
non-trivial braiding. For $n = 3$ one has
$\Conf_2(\bR^3) \simeq S^2 \times \bR^3$ (choosing the centre-of-mass
frame and the relative direction), hence $\pi_1(\Conf_2(\bR^3)) =
\pi_1(S^2) = 0$. Does this vanishing literally rule out
\emph{non-trivial braiding} on the $\Etwo$-restriction of an
$\Ethree$-algebra, or does the $\Etwo$-restriction still carry
non-trivial higher braiding (triple braiding, etc.)?
\end{attack}

\begin{heal}[The claim is \emph{symmetric} monoidal, not literally
trivial]
\label{heal:pi1}
The operadic restriction $\mathrm{res}_{3 \to 2}$ sends an
$\Ethree$-algebra~$B$ to the underlying $\Etwo$-algebra obtained
by forgetting the third disk direction. By Dunn additivity
$\Ethree \simeq \Etwo \otimes \Eone$, and the forgetful to
$\Etwo$ discards the $\Eone$ factor. The braiding on the
restriction is induced by $\pi_1(\Conf_2(\bR^3))$, which is
trivial. \emph{Triviality of $\pi_1$ means the braiding
$\beta\colon M \otimes N \to N \otimes M$ satisfies
$\beta_{N,M} \circ \beta_{M,N} = \id_{M \otimes N}$; equivalently
the braiding is \emph{symmetric}.} It does \emph{not} mean the
braiding is the literal permutation: a symmetric monoidal
$\Etwo$-algebra has a braiding but an involutive one.

Symbolically: the braiding factors through the symmetric group,
$\beta^2 = \id$. The restriction
$\mathrm{res}_{3 \to 2}\colon \Ethree\text{-alg} \to
\Etwo\text{-alg}$ factors through the full
subcategory $\Etwo^{\mathrm{sym}}\text{-alg} \subset
\Etwo\text{-alg}$ of symmetric-monoidal $\Etwo$-algebras, which
is equivalent to the category of $\Einf$-algebras. In particular
no $\bZ$-many braidings survive; exactly one involution
survives per pair of objects.

The higher braidings (triple, etc.) are controlled by
$\pi_n(\Conf_2(\bR^3)) = \pi_n(S^2)$, which is non-trivial for
$n = 2$ (a $\bZ$ Hopf class), $n \geq 3$ (finite). These higher
braidings do exist and contribute to the full $\Ethree$-structure.
What vanishes is specifically the $\pi_1$-level braiding, and that
is what the quantum-group $R$-matrix would require. The
$R$-matrix is a $\pi_1$-class (the half-monodromy along a loop
in $\Conf_2(\bC)$); its image in $\Conf_2(\bR^3)$ is zero.

The claim in the target is therefore to be read:
\textit{Direct restriction $\Ethree \to \Etwo$ produces only
symmetric braidings; in particular the non-trivial quantum-group
$R$-matrix cannot arise from this arrow.} That is exactly the
content of Vol~III
\texttt{chapters/theory/cy\_to\_chiral.tex:196} and Vol~III
\texttt{drinfeld\_center.tex:404}.
\end{heal}

\subsection*{Attack~2 --- $\Etwo$-chiral (boundary) vs
$\Eone$-chiral (bar/QG): both carry~$z$}

\begin{attack}[What distinguishes the $\Etwo$-chiral at the boundary
from the $\Eone$-chiral of the bar/QG?]
\label{att:e2-vs-e1}
Both objects live on a complex curve. Both carry a spectral-like
variable~$z$. Arrow~2 sends the former to the latter via the
ordered bar complex. Surely both are natively factorisation
algebras on $\Ran(C)$. What is the operadic content of the
passage?
\end{attack}

\begin{heal}[$\Etwo$-chiral has symmetric OPE and an
$\Einf$-factorisation; $\Eone$-chiral retains ordered collision
data]
\label{heal:e2-vs-e1}
The $\Etwo$-chiral boundary algebra $\cA$ arriving from arrow~1
is the factorisation algebra of observables on $C$ in the sense
of Costello--Gwilliam / Francis--Gaitsgory: a cosheaf on $\Ran(C)$
with symmetric-monoidal OPE, compatible with disjoint-union
factorisation on ordered \emph{and} unordered Ran. The
$z$-variable here is the complex coordinate on~$C$ and the
operad that acts is $\Etwo^{\mathrm{hol}}$ (notion (A) of
Def.~\ref{def:three-e2}).

Arrow~2, the ordered bar complex
$\barB^{\mathrm{ord}}(\cA) = T^c(s^{-1}\bar\cA)$ built on
$\Conf_n^{\mathrm{ord}}(C)$ / $\mathrm{FM}_n^{\mathrm{ord}}(C)$,
outputs an object on \emph{ordered} Ran. The ordered bar retains
the braid-group action on collision strata; the symmetric bar
loses it by projection to $\Sigma_n$-coinvariants. On
$\Einf$-chiral inputs (Heisenberg, $V_k(\fg)$, Virasoro) the
$\Sigma_n$-coinvariant projection is a quasi-isomorphism
$\barB^{\mathrm{ord}} \to \barB^\Sigma$ (Vol~I
Theorem~\texttt{thm:e1-formality-bridge}); on genuine
$\Eone$-chiral inputs (Yangian $Y^+(\widehat\fgl_1)$,
Etingof--Kazhdan quantum vertex, outputs of $\Phi_3$ at
CoHA-inputs) the projection is lossy with fibre encoding Arnold
braid-group representations.

\emph{Conceptually.} The distinction is not ``two different
$z$-variables'' but ``two different operadic structures on the same
underlying factorisation algebra.'' Boundary: full $\Sigma_n$
action survives (symmetric braiding, $\pi_1$-trivial restriction
from $\Ethree$). Bar: only the ordered braid action survives,
and the braid-to-symmetric projection map --- which is universally
the quotient by the Arnold braid-group $B_n \twoheadrightarrow
\Sigma_n$ --- is not a quasi-isomorphism when the algebra is
genuinely $\Eone$-chiral.

\emph{Functorial form.} Arrow~2 is the forgetful / restriction
along the map of operads $\Etwo^{\mathrm{hol}} \leftarrow
\Eone^{\mathrm{hol}}$, made explicit by the ordered-bar
construction applied to the underlying factorisation structure.
On $\Eone\text{-chiral}$ this is the Vol~I bar functor
$B^{\mathrm{ord}}$ of Theorem~A; on $\Einf$-chiral the two
collapse to the same object up to quasi-isomorphism.

At the level of $(\infty,1)$-categories, the adjunction is:
\[
\Omega^{\mathrm{ch}} \dashv B^{\mathrm{ord}}\colon
\Eone\text{-HolFA}(C)^{(\infty,1)}
\longrightarrow
\mathrm{ChirCoAlg}^{\mathrm{conil}}(C)^{(\infty,1)},
\]
Vol~III \texttt{chapters/theory/cy\_to\_chiral.tex:250}. The
source is $\Eone$-chiral; the target is the conilpotent chiral
coalgebras of Vol~I; the adjunction is the chiral bar--cobar
pair of Theorem~A, and it is $(\infty,1)$-functorial without
chain-level choices.
\end{heal}

\subsection*{Attack~3 --- closure: does arrow~5 land at the
\emph{same} $\Ethree$ as the source of arrow~1?}

\begin{attack}[Is the target $\Ethree^{\mathrm{top}}(\text{derived
centre})$ literally the source $\Ethree^{\mathrm{top}}(\text{bulk})$,
or just an $\Ethree$ of the same type?]
\label{att:closure}
The symbol $\Ethree^{\mathrm{top}}$ appears twice. At the source it
is the $3$d HT bulk observable algebra: a topological factorisation
algebra on $C \times \bR$ whose restriction to the defect curve
is the boundary $\cA$. At the target it is the chiral derived
centre $\Zder(\cA) = C^\bullet_{\mathrm{ch}}(\cA, \cA)$, with
$\Ethree$-structure via the higher Deligne conjecture. Both are
$\Ethree$-objects, but are they the same $\Ethree$-algebra?
\end{attack}

\begin{heal}[Closure is a conjecture; the two $\Ethree$'s are the
same algebra on the Koszul/topologisation locus]
\label{heal:closure}
The identification of the bulk with the derived centre is
Conjecture~\texttt{conj:v1-drinfeld-center-equals-bulk}
(Vol~I \texttt{chapters/frame/preface.tex:4525}):
\[
Z(U_\cA) \;\simeq\; \Zder(\cA)
\quad\text{as $\Ethree^{\mathrm{top}}$-algebras},
\]
where $U_\cA = \cA \bowtie \cA^!$ is the Drinfeld double and
$Z(-)$ is the Drinfeld centre. Three obstructions are identified
(Vol~I preface 10.6): pointwise-reduction failure for class~$M$,
Verdier-dual completion of $\cA^!$, and bar--cobar landing at
$\cA$ (not at its centre). The conjecture is open in general.

\emph{Partial closure.} The conjecture is proved on specific
strata:
\begin{itemize}
\item For simple~$\fg$, the $\Ethree^{\mathrm{top}}$-deformation
 space of $\Zder(V_k(\fg))$ is $\hbar H^3(\fg)[[\hbar]]$, which
 is one-dimensional, and formal-disk restriction recovers the
 Costello--Francis--Gwilliam perturbative Chern--Simons
 $\Ethree^{\mathrm{top}}$-algebra
 (Vol~I \texttt{chapters/frame/preface.tex:4554}).
\item For the CY$_3$ input $\cC = \CoHA(\bC^3)$,
 $\Phi_3(\cC) = Y^+(\widehat\fgl_1)$ and
 $\cZ(\Rep^{\Eone}(Y^+(\widehat\fgl_1))) \simeq
 \Rep^{\Etwo}(Y(\widehat\fgl_1)) \simeq
 \Rep^{\Etwo}(\cW_{1+\infty})$ at $c=1$
 (Vol~III Theorem~\texttt{thm:c3-drinfeld-center-restate},
 \texttt{drinfeld\_center.tex:588}).
\end{itemize}

The name ``circle'' is hence a \emph{conjectural closure}; the
five arrows are all proved unconditionally or under specified
hypotheses, but arrow~5 lands at an $\Ethree^{\mathrm{top}}$-object
that is only conjecturally equal to the source
$\Ethree^{\mathrm{top}}$. The honest label is \emph{spiral}; the
precise form is: the five arrows form a zig-zag in
$(\infty, 1)$-operadic data whose closure as a literal circle is
conj.~\texttt{conj:v1-drinfeld-center-equals-bulk}. The current
programme commits that the closure is available on the
Koszul--topologisation locus (class~$\mathbf{L}$ at non-critical
level, plus the $\bC^3$ class~$\mathbf{C}$ patch) and conjectural
elsewhere.

\emph{Refinement.} The two $\Ethree^{\mathrm{top}}$'s are not
the same \emph{factorisation-algebra kind}: the bulk is a
topological factorisation algebra on $C \times \bR$; the derived
centre is the algebraic Hochschild object with its Deligne
$\Ethree$-structure. The identification requires (i) the
topologisation of the $\SCchtop$-algebra arising from $\cA$ to
an $\Ethree^{\mathrm{top}}$-object (proved at non-critical
Kac--Moody; conjectural for general Virasoro / $\cW_N$), and
(ii) the translation between Costello--Gwilliam HT
factorisation algebras and algebraic Hochschild complexes
(Vol~II \texttt{chapters/theory/chiral\_higher\_deligne.tex:956}).
\end{heal}

\subsection*{Attack~4 --- Drinfeld centre as \emph{right} adjoint:
is that the correct $(\infty,n)$-framework?}

\begin{attack}[Is $\cZ \dashv U$ the right $(\infty,1)$-statement, or
is the correct framework $(\infty, 2)$-categorical?]
\label{att:drinfeld-adjoint}
The Drinfeld centre is presented as the right adjoint to the
forgetful $U\colon \Etwo\text{-Cat} \to \Eone\text{-Cat}$. But
the passage $\Eone \to \Etwo$ on categories is naturally
$(\infty, 2)$-categorical (braided monoidal categories are
$3$-categorical objects by the cobordism hypothesis). Is the
adjunction in the correct homotopy universe?
\end{attack}

\begin{heal}[Both readings apply; the $(\infty,1)$-statement is Lurie
\emph{HA}~5.3.1.30, the $(\infty,2)$-statement is the
BZFN theorem]
\label{heal:drinfeld-adjoint}
Two adjunction statements coexist at different homotopy depths.

\emph{Level $(\infty, 1)$.} Lurie, \emph{Higher Algebra}
5.3.1.30, constructs the Drinfeld centre of an
$\Eone$-monoidal $(\infty,1)$-category as the right adjoint
to the forgetful from $\Etwo$-monoidal $(\infty,1)$-categories
to $\Eone$-monoidal $(\infty,1)$-categories. Concretely, for
$\cC$ a monoidal $(\infty, 1)$-category,
\[
\cZ(\cC) \;=\;
\mathrm{Fun}^{\cC\text{-lin}}_{\mathrm{mon}}(\cC, \cC),
\]
the $(\infty, 1)$-category of $\cC$-linear monoidal endofunctors,
carries a canonical $\Etwo$-monoidal (equivalently, braided
monoidal) structure, and the functor $\cZ$ right-adjoins the
forgetful.

\emph{Level $(\infty, 2)$.} Ben-Zvi--Francis--Nadler
\cite{BFN2010} identify, for $\cC = \Rep^{\Eone}(A)$ with $A$ an
$\Eone$-algebra in a symmetric monoidal stable
$(\infty, 1)$-category~$\cS$:
\[
\cZ(\Rep^{\Eone}(A))
\;\simeq\;
\Rep^{\Etwo}\bigl(\HH^\bullet(A, A)\bigr)
\quad\text{in $\Etwo\text{-Cat}^{(\infty,1)}$}.
\]
This is an equivalence of $\Etwo$-monoidal $(\infty,1)$-categories,
i.e., an equivalence of $3$-categorical / $(\infty,2)$-categorical
braided monoidal data. Both sides describe the same braided
monoidal structure; the BZFN theorem is the categorified
Hochschild-Deligne identification. Vol~III
\texttt{chapters/theory/drinfeld\_center.tex:101--114} states it
explicitly:
\[
\cZ(\Rep^{\Eone}(A)) \;\simeq\; \Rep^{\Etwo}(\Zder(A)),
\qquad
\Zder(A) = \RHom_{A\otimes A^{\op}}(A, A).
\]

\emph{Confirmation of the right-adjoint reading.} The
forgetful $U\colon \Etwo\text{-Cat} \to \Eone\text{-Cat}$ has a
right adjoint (Lurie 5.3.1.30) on presentable $(\infty, 1)$-categories,
and the right adjoint sends $\cC \mapsto \cZ(\cC)$. At
$\infty$-categorical level the adjunction is \emph{universal}:
every $\Etwo$-lift of an $\Eone$-monoidal $\cC$ factors uniquely
through $\cZ(\cC)$. On coherent $\Etwo$-categories coming from
$\Etwo$-algebras, the BZFN identification then replaces the
right-adjoint construction with an algebraic one in terms of
$\HH^\bullet$ (Deligne conjecture $\Rightarrow$ $\Etwo$ on
$\HH^\bullet(A, A)$).

\emph{Universal property in one sentence.} $\cZ(\cC)$ is the
universal braided monoidal category equipped with a
$\cC$-linear monoidal forgetful functor; equivalently, it is the
endomorphism $\Etwo$-category of the $\Eone$-module $\cC$
acting on itself.

The circle's arrow~3 is therefore the \emph{right adjoint} to the
forgetful; it is
NOT the averaging map (which is the left-type
$\Sigma_n$-coinvariant projection, losing ordered data), and it
is NOT arrow~2 going backward (which would require $R$-matrix
data and be circular). The Drinfeld-centre construction creates
$\Etwo$-braided data out of $\Eone$-monoidal data by adjoining
half-braidings; the BZFN theorem identifies this data with the
$\Etwo$-algebra structure on the Hochschild cochains produced by
the classical Deligne conjecture.
\end{heal}

\subsection*{Attack~5 --- could the circle close via \emph{two}
volumes? Or split across \emph{four}?}

\begin{attack}[The three-volume decomposition is claimed \emph{forced}.
Test by considering two-volume and four-volume alternatives]
\label{att:volumes}
\textbf{Two volumes.} Suppose we merge the Vol~II HT-boundary story
into Vol~I. The five arrows would then live in: Vol~I (all five)
and Vol~III (CY source). Does this actually close? What provides
the geometric source of $\Eone$-chirality?
\textbf{Four volumes.} Suppose we split Vol~II along the seam
between the $\SCchtop$ content (Swiss-cheese chain-level) and the
$\Ethree$-topologisation via Sugawara (when it applies). Does this
split produce two distinct mathematical programmes?
\end{attack}

\begin{heal}[Two volumes fail because there is no geometric source
of $\Eone$-chirality without CY$_d$, $d \geq 3$; four volumes
split along a \emph{fake} seam because $\SCchtop$ and $\Ethree$-topologisation
are two faces of the same intermediary]
\label{heal:volumes}

\emph{Why two volumes fail.} The five-arrow circle starts at the
bulk $\Ethree^{\mathrm{top}}$ observable algebra and must produce,
as arrow~1's output, a genuine $\Eone$-chiral algebra at the
defect curve. On the $\Einf$-chiral subcategory (standard VA
landscape: Heisenberg, $V_k(\fg)$, Virasoro, $\cW_N$), the
ordered bar $\barB^{\mathrm{ord}}$ coincides quasi-isomorphically
with the symmetric bar $\barB^\Sigma$, and the arrow~3 Drinfeld
centre on $\Rep^{\Einf}(\cA) = \Rep^{\Einf}(\cA)$ collapses to the
identity (the Drinfeld centre of a symmetric monoidal category is
itself, Vol~III \texttt{drinfeld\_center.tex:2271}). In that
regime the circle is degenerate.

The genuinely $\Eone$-chiral algebras --- Yangian
$Y^+(\widehat\fgl_1)$, Etingof--Kazhdan quantum vertex,
$\Eone$-chiral shadows of CoHAs of CY$_3$-varieties, K3-Yangian
of $D^b(\Coh(K3 \times E))$ --- do not live on curves alone.
They arise as Stage-$2$ specialisations of Stage-$1$
$E_d$-holomorphic factorisation algebras on CY$_d$-varieties
for $d \geq 3$
(Vol~III \texttt{chapters/theory/cy\_to\_chiral.tex:30--70}):
\[
\Phi_d \;=\; \mathrm{Sp}^{\mathrm{ch}}_{\Sigma_{d-1}, C}
\circ \Phi^{\mathrm{FA}}_d.
\]
This factorisation forces a CY$_d$ input for $d \geq 3$; neither
the curve-based chiral geometry of Vol~I nor the $3$d HT slab of
Vol~II provides this input. \textbf{Collapsing Vol~III into Vol~I
or Vol~II removes the geometric source of
$\Eone$-chirality}, and the circle has no non-trivial instances:
every example reduces to the $\Einf$-chiral landscape and arrow~3
trivialises. Concretely: without Vol~III, the programme cannot
produce the Yangian of $\widehat\fgl_1$ as a geometric shadow, and
the Drinfeld-centre arrow is empty.

\emph{Why four volumes split along a fake seam.} The proposal is
to separate $\SCchtop$ (the Swiss-cheese two-coloured operad,
working at chain-level, with no Sugawara assumption) from the
$\Ethree$-topologisation (requiring an inner conformal vector~$\nu$
and the Sugawara identity
$T_{\mathrm{Sug}} = [Q, G]$; proved at non-critical Kac--Moody,
conjectural elsewhere).

Inside Vol~II these are not two arguments but two layers of one
argument. The $\SCchtop$-algebra is the \emph{chain-level}
(boundary + bulk) structure that an $\Eone$-chiral algebra
manifests when promoted to a two-coloured operad with
directionality constraint (open cannot feed closed,
Vol~I \texttt{chapters/frame/preface.tex:4395}). Dunn additivity
does not apply to $\SCchtop$, so $\SCchtop \neq \Ethree$ on the
nose. The topologisation theorem
(Vol~II \texttt{V2-thm:iterated-sugawara-construction} and Vol~I
\texttt{thm:topologization}) adds an inner conformal vector
that \emph{identifies} the $\SCchtop$-algebra with an
$\Ethree^{\mathrm{top}}$-algebra, provably at non-critical
affine Kac--Moody (Sugawara closed-form), cohomologically
general, chain-level conjectural for class~$\mathbf{M}$
(Virasoro, $\cW_N$) on the original complex.

Splitting this into two volumes would require presenting
$\SCchtop$ as a freestanding theory disjoint from its
Sugawara-topologisation; but the sole reason the programme needs
$\SCchtop$ at all is to have a \emph{chain-level} intermediary
between $\Eone$-chiral and $\Ethree^{\mathrm{top}}$ within which
topologisation can be proved. Separating $\SCchtop$ from the
Sugawara argument loses the proof of topologisation, because
topologisation is literally the identity that upgrades
$\SCchtop$ to $\Ethree^{\mathrm{top}}$ when $\nu$ exists.

\emph{Conclusion.} Three volumes are precisely the minimal
decomposition that (a) carries the five arrows as
algebraic-geometric theorems (Vol~I); (b) provides the
chain-level / BRST / 3d HT physics and the $\SCchtop$
topologisation (Vol~II); (c) supplies the CY$_d$ geometric
source of $\Eone$-chirality (Vol~III). Two collapse (b) or (c)
away; four severs (b) along a non-existent seam. This is the
precise sense in which the circle \emph{forces} three volumes.
\end{heal}

\section{Rigorous step-by-step account of the five arrows}

Each arrow specified as an $(\infty, 1)$-functor, with witnessing
algebras and explicit categorical input / output.

\subsection*{Arrow~1. $\Ethree^{\mathrm{top}}(\mathrm{bulk}) \to
\Etwo\text{-chiral}(\mathrm{bdy})$: bulk-to-defect restriction}

\begin{defn}[Arrow~1 functor]
Let $B$ be an $\Ethree^{\mathrm{top}}$-algebra of observables of a
$3$d HT theory on the slab $C \times \bR_{\geq 0}$ (Vol~II
\texttt{chapters/theory/foundations\_recast\_draft.tex:549}).
Restriction to the codimension-$2$ defect $C \times \{0\}$
yields
\[
\mathrm{res}\colon\;
\Ethree^{\mathrm{top}}\text{-alg}(C \times \bR_{\geq 0})
\;\longrightarrow\;
\Etwo\text{-alg}^{\mathrm{hol}}(C),
\qquad
B \longmapsto \cA := B|_{C \times \{0\}}.
\]
The operadic-level drop by one comes from the codimension-$2$
embedding: the rank-$2$ normal bundle supplies the two extra
disk-directions lost on the defect. In coordinates:
$\mathrm{res} = \pi_*(j^* B)$ where $j\colon C \hookrightarrow
C \times \bR_{\geq 0}$ is the inclusion and $\pi\colon C \to C$
is the identity; the holomorphic-factorisation structure on
$\cA$ survives the restriction because $B$ is topological in the
transverse $\bR$-direction.
\end{defn}

\begin{rem}[Witnessing algebras]
\begin{itemize}
\item $B = $ $3$d abelian Chern--Simons at level $k$ on
$C \times \bR_{\geq 0}$; defect restriction yields
$\cA = H_k$, the Heisenberg chiral algebra at level $k$.
Class $\mathbf{G}$.
\item $B = $ $3$d non-abelian Chern--Simons at level $k$ with
gauge group $G$ (Lie algebra $\fg$) and $k \neq -h^\vee$;
defect yields $\cA = V_k(\fg)$, the affine vertex algebra at
level $k$. Class $\mathbf{L}$.
\item $B = $ $3$d $\beta\gamma$-system (holomorphic twist of
$\cN = 2$ $\mathrm{U}(1)$); defect yields
$\cA = \beta\gamma(\lambda)$ at weight $\lambda$. Class
$\mathbf{C}$.
\item $B = $ $3$d HT boundary of $(2, 0)$ theory with conformal
vector; defect yields $\cA = \mathrm{Vir}_c$ or $\cW_N$ at
matching central charge. Class $\mathbf{M}$.
\item $B = $ $3$d HT observables dual to K3 $\times E$ via
BPS geometry; defect yields the K3-Yangian chiral algebra
$\mathbf{H}_{\Delta_5}$. Cross-volume class $\mathbf{B}$
(Vol~III).
\end{itemize}
Each restriction is a theorem (for class $\mathbf{G, L, C}$
unconditional at non-critical level), with full chain-level
control when the boundary condition is topologised.
\end{rem}

\subsection*{Arrow~2. $\Etwo\text{-chiral}(\mathrm{bdy}) \to
\Eone\text{-chiral}(\mathrm{bar/QG})$: ordered bar complex}

\begin{defn}[Arrow~2 functor]
The ordered bar functor on $\Etwo$-chiral (equivalently,
$\Einf$-chiral when the input is a standard VA) algebras:
\[
\barB^{\mathrm{ord}}\colon
\Eone\text{-HolFA}(C)^{(\infty,1)}
\;\longrightarrow\;
\mathrm{ChirCoAlg}^{\mathrm{conil}}(C)^{(\infty,1)},
\qquad
\cA \longmapsto T^c(s^{-1}\bar\cA),
\]
with chiral bar differential
$d_\cB = d_{\mathrm{int}} + d_{\mathrm{mult}}$ built from
collision residues on $\mathrm{FM}^{\mathrm{ord}}_n(C)$
(Vol~I Theorem~\texttt{thm:bar-cobar-isomorphism-main}).
Right-adjoint to chiral cobar $\Omega^{\mathrm{ch}}$:
$\Omega^{\mathrm{ch}} \dashv \barB^{\mathrm{ord}}$.
\end{defn}

\begin{rem}[Witnessing]
For $\cA = V_k(\fg)$, $\barB^{\mathrm{ord}}(V_k(\fg))$ is the
$\Eone$-chiral coalgebra whose primitive filtration carries the
universal $R$-matrix $r^{\mathrm{KM}}(z) = k\,\Omega/z$ of the
affine Kac--Moody Yangian. For $\cA = H_k$,
$\barB^{\mathrm{ord}}(H_k)$ is the chiral Hopf coalgebra of the
Heisenberg Yangian, with $r^{\mathrm{Heis}}(z) = k/z$. For
$\cA = Y^+(\widehat\fgl_1)$ (arising at $d = 3$ from
$\Phi_3(\CoHA(\bC^3))$), $\barB^{\mathrm{ord}}(\cA)$ is the bar
coalgebra of the positive half-Yangian; its Drinfeld double
(equivalently, the Verdier dual $\cA^!$, equivalently the chiral
derived centre) is the full $Y(\widehat\fgl_1)$ which, at $c=1$,
recovers $\cW_{1+\infty}$.
\end{rem}

\subsection*{Arrow~3. $\Eone\text{-chiral} \to \Etwo\text{-braided}$:
Drinfeld centre of the monoidal representation category}

\begin{defn}[Arrow~3 functor]
For $\cA$ an $\Eone$-chiral algebra, the Drinfeld centre of its
chiral representation category:
\[
\cZ\bigl(\Rep^{\Eone}(\cA)\bigr)
\;\in\;
\Etwo\text{-BrMon}\text{-Cat}^{(\infty,1)}.
\]
Objects are pairs $(M, \sigma_M)$ with $M \in \Rep^{\Eone}(\cA)$
and $\sigma_M\colon M \otimes (-) \xrightarrow{\sim} (-) \otimes M$
a half-braiding satisfying the hexagon axiom. BZFN theorem
identifies:
\[
\cZ(\Rep^{\Eone}(\cA)) \;\simeq\; \Rep^{\Etwo}(\Zder(\cA))
\quad\text{in $\Etwo\text{-BrMon}\text{-Cat}^{(\infty,1)}$},
\]
$\Zder(\cA) = \RHom_{\cA\otimes \cA^{\op}}(\cA, \cA) =
\HH^\bullet(\cA, \cA)$
(Vol~III \texttt{chapters/theory/drinfeld\_center.tex:108--113}).
\end{defn}

\begin{rem}[Witnessing: the Heisenberg half-braiding]
For $\cA = H_k$, the Fock modules $\cF_\lambda$ carry a
half-braiding
$\sigma_{\cF_\lambda}(\cF_\mu) = e^{k\lambda\mu/z}$
(the normal-ordered exponential of the Heisenberg OPE). The
classical $r$-matrix $r^{\mathrm{Heis}}(z) = k\,\Omega/z$ is the
first-order term of $\log R(z)$; the averaging map annihilates
the full $R$-matrix to the scalar $\kappa_{\mathrm{ch}} = k$.
Vol~III \texttt{drinfeld\_center.tex:547--579}.

For $\cA = Y^+(\widehat\fgl_1)$, BZFN gives
\[
\cZ(\Rep^{\Eone}(Y^+(\widehat\fgl_1))) \;\simeq\;
\Rep^{\Etwo}(Y(\widehat\fgl_1)) \;\simeq\;
\Rep^{\Etwo}(\cW_{1+\infty}|_{c = 1}),
\]
with the middle equivalence by Drinfeld doubling
$Y(\widehat\fgl_1) = Y^+(\widehat\fgl_1) \bowtie
(Y^+(\widehat\fgl_1))^{\op}$ and the right equivalence by
Procházka--Rapčák's vertex algebra isomorphism at $c = 1$.
Vol~III \texttt{drinfeld\_center.tex:588--608}.
\end{rem}

\subsection*{Arrow~4. $\Etwo\text{-braided} \to
\Ethree^{\mathrm{top}}$: higher Deligne conjecture}

\begin{defn}[Arrow~4 functor]
For $B$ an $\Etwo$-algebra (equivalently, a braided monoidal
$(\infty, 1)$-category coming from an $\Etwo$-algebra via
$\Rep^{\Etwo}$), the Hochschild cochain complex
$\HH^\bullet(B, B) = \RHom_{B \otimes_{\Etwo} B^{\op}}(B, B)$
carries an $\Ethree$-structure (Lurie, \emph{Higher Algebra}
Theorem~5.3.1.30). On the $(\infty, 1)$-category of
$\Etwo$-monoidal $(\infty, 1)$-categories:
\[
\HH^\bullet_{\mathrm{cat}}\colon
\Etwo\text{-Cat}^{(\infty,1)}
\;\longrightarrow\;
\Ethree\text{-Cat}^{(\infty,1)}.
\]
\end{defn}

The passage is the higher Deligne theorem. The subtle point is
that $\Ethree^{\mathrm{top}}$ (full topological $E_3$) requires the
Sugawara conformal vector (Vol~I
\texttt{thm:topologization});
the bare $\Ethree$-structure produced by Deligne carries two
real disk-directions and one direction whose topologisation
requires $\nu$. For affine Kac--Moody at non-critical level
Sugawara closes this on the nose; for class~$\mathbf{M}$ the
chain-level $\Ethree^{\mathrm{top}}$ is conjectural on the
original complex.

\subsection*{Arrow~5 (closure). $\Ethree^{\mathrm{top}}(\text{der.
centre}) \xrightarrow{\,\text{conj.}\,}
\Ethree^{\mathrm{top}}(\text{bulk})$}

\begin{defn}[The closure conjecture]
Conjecture
\texttt{conj:v1-drinfeld-center-equals-bulk}
\textup{(}Vol~I \texttt{chapters/frame/preface.tex:4525}\textup{):}
\[
Z(U_\cA) \;\simeq\; \Zder(\cA) \quad
\text{in $\Ethree^{\mathrm{top}}\text{-alg}$}.
\]
Proved on: simple $\fg$ (Whitehead rigidity), $\fgl_N$ (explicit
$\mathrm{GL}(N)$-pinning), semisimple, reductive, nilpotent (via
Heisenberg $\dim = 4$ and $\mathrm{Sym}^2(\mathfrak{h}^*)^W$),
solvable (Whitehead failure explicit, solvable Lie obstruction
finite). Chain-level on qi-model via
Fresse Vol~II Thm 16.2.1. Three obstructions remain open in
general: (a) pointwise reduction for class~$\mathbf{M}$; (b)
Verdier-dual factorisation of $\cA^!$; (c) bar--cobar
identification with the centre via Hochschild. Resolution of
any one reduces the conjecture to $E_2$-comparison, which
follows from higher Deligne.
\end{defn}

\section{Concrete verification: $\cZ(\Rep^{\Eone}(Y^+(\widehat\fgl_1)))
\simeq \Rep^{\Etwo}(\cW_{1+\infty}|_{c=1})$}

The key concrete witness. Direct computation (Vol~III
\texttt{drinfeld\_center.tex:582--608}):

\begin{prop}[$\bC^3$ Drinfeld-centre identification]
\label{prop:c3-verify}
Let $\cA = Y^+(\widehat\fgl_1)$ be the positive half of the
affine Yangian of $\widehat\fgl_1$, arising as
$\Phi_3(\CoHA(\bC^3))$ at Stage-$2$ specialisation
$(\Sigma_2, C) = (\mathrm{pt}, \bC)$. Then
\[
\cZ(\Rep^{\Eone}(\cA)) \;\simeq\; \Rep^{\Etwo}(Y(\widehat\fgl_1))
\;\simeq\; \Rep^{\Etwo}(\cW_{1+\infty})|_{c=1}
\]
as $\Etwo$-braided monoidal $(\infty, 1)$-categories.
\end{prop}

\begin{proof}[Sketch; attributed]
The first equivalence is the Drinfeld-doubling identity
$Y(\widehat\fgl_1) = Y^+(\widehat\fgl_1) \bowtie
(Y^+(\widehat\fgl_1))^{\op}$ combined with BZFN
(Ben-Zvi--Francis--Nadler 2010). The second is the vertex
algebra isomorphism $Y(\widehat\fgl_1)\bigr|_{c=1} \simeq
\cW_{1+\infty}\bigr|_{c=1}$ (Procházka--Rapčák 2019 / Litvinov
2020). Evaluation modules $V_u$ of $Y(\widehat\fgl_1)$
carry a spectral parameter $u \in \bC$, and the half-braiding
$\sigma_{V_u}(V_v) = R(z)$ with $z = u - v$ specialises to the
rational $R$-matrix of the Yangian, equivalently to the
$\cW_{1+\infty}$ $R$-matrix of Maulik--Okounkov at $c = 1$.
\end{proof}

\emph{Check.} On the side $\cZ(\Rep^{\Eone}(Y^+(\widehat\fgl_1)))$
the $\Etwo$-braiding datum is the half-braiding universal
adjunction (Step~3 of \texttt{drinfeld\_center.tex:422--441}).
On the side $\Rep^{\Etwo}(\cW_{1+\infty})$ the $\Etwo$-braiding
is the native braided monoidal structure of representations of
the $\cW_{1+\infty}$ vertex operator algebra at central
charge~$1$ (a genuine $\Etwo$-structure because
$\cW_{1+\infty}$ is a standard chiral algebra,
$\Einf$-chiral, and $\Rep$-categories of $\Einf$-chiral algebras
inherit a modular tensor category structure when the algebra
is rational). The two sides agree by $R$-matrix matching:
the half-braiding $\sigma_{V_u}(V_v) = e^{k \lambda \mu / z}$
on Fock-sector Heisenberg restriction, $\widetilde R(z)$ on
generic evaluation, rational $R$-matrix with pole structure
$R(z) = 1 + \hbar r(z) + O(\hbar^2)$ with
$r^{\mathrm{Heis}}(z) = k\,\Omega/z$. All three presentations
match on the Heisenberg sub-Yangian, and the full comparison
for $Y(\widehat\fgl_1)$ is Maulik--Okounkov's explicit
$R$-matrix formula.

Vol~III \texttt{drinfeld\_center.tex:360--374} records a
$98$-test compute verification
(\texttt{compute/lib/derived\_vs\_drinfeld\_infty.py}) covering
the K3 dimensions ($325$ vs $49$, Euler characteristic $277$),
$\Etwo$ matching across shadow classes, and TCFT tower
structure $\{b, B^{(2)}\} = 0$.

\section{Summary}

The five-arrow circle $\Ethree \to \Etwo \to \Eone \to \Etwo \to
\Ethree$ has the following honest status and structure:

\begin{itemize}
\item Arrow~1 (restriction to defect): proved for
$\Einf$-chiral class $\mathbf{G, L, C}$ at non-critical level;
inscribed in Vol~II locality and axioms chapters.

\item Arrow~2 (ordered bar): proved unconditionally for all
$\Eone$-chiral algebras with conilpotent coalgebra bar;
adjunction $\Omega^{\mathrm{ch}} \dashv \barB^{\mathrm{ord}}$
(Vol~I Theorem~A).

\item Arrow~3 (Drinfeld centre): proved unconditionally as the
right adjoint to forgetful in Lurie \emph{HA}~5.3.1.30; the
BZFN theorem identifies it with $\Rep^{\Etwo}(\HH^\bullet(A, A))$
for $A$ an $\Eone$-algebra in a presentable stable
$(\infty,1)$-category. Verified concretely for the $\bC^3$
Yangian via Schiffmann--Vasserot / Maulik--Okounkov /
Procházka--Rapčák.

\item Arrow~4 (higher Deligne): proved unconditionally (Lurie,
Francis); $\HH^\bullet$ of an $\Etwo$-algebra carries
$\Ethree$-structure.

\item Arrow~5 (closure): conjectural in general, proved on
Koszul--topologisation locus (affine Kac--Moody non-critical;
$\bC^3$ CoHA patch); three obstructions identified (Vol~I
preface~10.6).
\end{itemize}

The $\pi_1(\Conf_2(\bR^3)) = 0$ fact rules out \emph{non-symmetric}
braiding from direct $\Ethree \to \Etwo$ restriction (the image is
$\Einf$-type / symmetric monoidal). The Drinfeld-centre arrow~3 is
therefore the unique source of non-trivial braiding in the circle.
Three volumes are forced: Vol~I carries the five arrows
algebraic-geometrically; Vol~II carries the $3$d HT /
$\SCchtop$ / topologisation chain-level content; Vol~III provides
the CY$_d$ geometric source of $\Eone$-chirality without which
the circle is empty.

\paragraph{Attribution.} Chapters read:
\texttt{/Users/raeez/chiral-bar-cobar/chapters/frame/preface.tex},
\texttt{chapters/theory/en\_koszul\_duality.tex},
\texttt{chapters/frame/part\_iv\_platonic\_introduction.tex};
\texttt{/Users/raeez/calabi-yau-quantum-groups/chapters/theory/cy\_to\_chiral.tex},
\texttt{chapters/theory/drinfeld\_center.tex};
plus Lurie \emph{Higher Algebra}~5.3.1.30, Ben-Zvi--Francis--Nadler
\emph{Integral Transforms} 2010, Schiffmann--Vasserot 2013,
Maulik--Okounkov 2019, Procházka--Rapčák 2019.

\end{document}
